\documentclass[12pt,a4paper]{article}
\usepackage[margin=25mm, headheight=15pt, headsep=10pt]{geometry}
\usepackage{fontspec}
\usepackage[table]{xcolor} % Note: table option is required for rowcolors
\usepackage{titlesec}
\usepackage{tocloft}
\usepackage{fancyhdr}
\usepackage{booktabs}
\usepackage{tcolorbox}
\tcbuselibrary{skins, breakable}
\usepackage[colorlinks=true, linkcolor=emerald, urlcolor=emerald, bookmarks=true]{hyperref}

\definecolor{emerald}{RGB}{0,102,51}
\definecolor{darkred}{RGB}{139,0,0}
\definecolor{lightgray}{RGB}{245,245,245}

\usepackage[nil]{babel}
\babelprovide[import, main]{bengali}
\babelprovide[import]{arabic}
\babelfont{rm}[Renderer=HarfBuzz,Scale=1.0]{Noto Serif Bengali}
\babelfont{sf}[Renderer=HarfBuzz]{Noto Sans Bengali}
\babelfont{tt}[Renderer=HarfBuzz]{Noto Sans Bengali}
\babelfont[arabic]{rm}[Renderer=HarfBuzz,Scale=1.1]{Noto Naskh Arabic}

% Section styling
\titleformat{\section}{\Large\bfseries\color{emerald}}{\thesection.}{0.5em}{}
\titleformat{\subsection}{\large\bfseries\color{emerald!85!black}}{\thesubsection.}{0.5em}{}
\titleformat{\subsubsection}{\normalsize\bfseries\color{emerald!70!black}}{\thesubsubsection.}{0.5em}{}
\titlespacing*{\section}{0pt}{18pt}{8pt}

% Fancy Header/Footer setup
\pagestyle{fancy}
\fancyhf{} % clear all header and footer fields
\fancyhead[L]{\color{emerald}\small\textbf{\nouppercase{\leftmark}}}
\fancyfoot[L]{\scriptsize\color{black!50}{\lat{AI}}-সংকলিত, আলেম-যাচাইকৃত নয়}
\fancyfoot[C]{\color{emerald!70!black}\thepage}
\renewcommand{\headrulewidth}{0.4pt}
\renewcommand{\headrule}{\hbox to\headwidth{\color{emerald}\leaders\hrule height \headrulewidth\hfill}}
\renewcommand{\footrulewidth}{0pt}

% Custom boxes
% 1. Ayat/Hadith Box
\newtcolorbox{ayatbox}[1][]{
    enhanced, breakable, 
    colback=emerald!5, 
    colframe=emerald, 
    boxrule=0.5pt, 
    arc=3pt, 
    left=10pt, right=10pt, top=10pt, bottom=10pt,
    #1
}

% 2. Quote/Translation Box
\newtcolorbox{quotebox}[1][]{
    enhanced, breakable,
    frame hidden,
    colback=lightgray,
    borderline west={3pt}{0pt}{emerald},
    left=15pt, right=10pt, top=10pt, bottom=10pt,
    sharp corners,
    #1
}

% 3. AI-generated notice box
\newtcolorbox{warnbox}[1][]{
    enhanced, breakable,
    colback=red!4,
    colframe=darkred,
    boxrule=0.8pt,
    arc=3pt,
    left=10pt, right=10pt, top=10pt, bottom=10pt,
    #1
}

% Commands
\newcommand{\arab}[1]{\foreignlanguage{arabic}{#1}}
\newcommand{\kib}[1]{\textcolor{emerald}{\textbf{#1}}}

% Latin (English) fallback: Noto Serif Bengali has NO Latin glyphs
\newfontfamily\latfont[Renderer=HarfBuzz]{Noto Serif}
\newcommand{\lat}[1]{{\latfont #1}}

\setlength{\parskip}{5pt}
\setlength{\parindent}{0pt}

% Fix for missing bullet in Noto Serif Bengali
\renewcommand{\labelitemi}{$\bullet$}



\begin{document}
\begin{center}{\Large\bfseries\color{emerald} আল্লাহর রহমত ও তাওবা — \lat{09}-রমজানে-স্ত্রী-সহবাসকারী-সাহাবি}\end{center}
\vspace{1em}
\section*{ঘটনা ৯ — রমজানে স্ত্রী-সহবাসকারী সাহাবি: পাপ স্বীকার, নবীজির (সা.) স্নেহ ও নতুন শুরু}

\begin{warnbox}
 \textbf{গুরুত্বপূর্ণ নোটিশ (\lat{Important} \lat{Notice}):} এই কন্টেন্টটি \lat{AI} (কৃত্রিম বুদ্ধিমত্তা) দিয়ে প্রস্তুত ও সংকলিত। \lat{AI} কঠোর সূত্র-মান নীতি মেনে শুধু নির্ভরযোগ্য উৎস থেকে তথ্য সংগ্রহ করেছে — কেবল সহীহ/হাসান হাদিস ও সহীহ সনদের আসার রাখা হয়েছে, দুর্বল সনদ বাদ দেওয়া হয়েছে। তবে এটি কোনো আলেম বা মুহাদ্দিস কর্তৃক যাচাইকৃত নয়।
\end{warnbox}

\begin{quote}
\textbf{সম্পর্কিত ফাইল:} হাদিস ফাইল — পাপ স্বীকার ও ক্ষমা; ব্যবহারিক গাইড — তাওবার শর্ত।
\end{quote}

\medskip\hrule\medskip

\subsection*{১. সূত্র কার্ড (\lat{Source} \lat{Card})}

\begin{table}[h]\centering\small
\begin{tabular}{ll}
বিষয় & বিবরণ \\
\hline
\textbf{ঘটনা} & এক সাহাবি রমজানের রোজা রেখে স্ত্রীর সাথে সহবাস করে ফেললেন; নবীজির (সা.) কাছে পাপ স্বীকার করলেন; নবীজি (সা.) কাফফারার (\lat{expiation}) ব্যবস্থা বললেন এবং তাকে সান্ত্বনা দিলেন \\
\textbf{বর্ণনাকারী} & আবু হুরাইরাহ (রাঃ) \\
\textbf{গ্রন্থ ও নম্বর} & সহীহ বুখারি ১৯৩৬; সহীহ মুসলিম ১১১১ \\
\textbf{অধ্যায়} & বুখারি: কিতাবুস সাওম — \lat{“}রোজা রেখে সহবাসকারীর কাফফারা\lat{”}; মুসলিম: কিতাবুস সিয়াম — \lat{“}রমজানের দিনে সহবাস কঠোরভাবে হারাম ও কাফফারা ওয়াজিব\lat{”} \\
\textbf{সনদের মান} & \textbf{সহীহ} (বুখারি-মুসলিম) \\
\textbf{স্থান ও সময়} & মদিনা — নবীজির (সা.) জীবদ্দশায় \\
\hline
\end{tabular}
\end{table}

\medskip\hrule\medskip

\subsection*{২. সংক্ষিপ্ত পটভূমি (\lat{Background})}

রোজার দিনে স্ত্রীর সাথে সহবাস করা বড় পাপ (\lat{major} \lat{sin}) — কুরআন ও সুন্নাহয় স্পষ্ট নিষেধ। এই সাহাবি পাপটি করে ফেললেন; কিন্তু তার প্রতিক্রিয়াটি ছিল শিক্ষণীয় — তিনি \textbf{লুকালেন না}, বরং ভেঙে পড়া অবস্থায় নবীজির (সা.) কাছে এসে স্বীকার করলেন। নবীজি (সা.) তাকে ধমক দিলেন না — তাকে পাপ থেকে বের হওয়ার পথ (কাফফারা) দেখালেন এবং সবশেষে স্নেহ ও হাসিমুখে বিদায় দিলেন। এই ঘটনা থেকে বোঝা যায় — \textbf{পাপ স্বীকার করলে আল্লাহর রাসূলের (সা.) কাছ থেকে পাওয়া যায় পথনির্দেশনা (\lat{guidance}) ও সান্ত্বনা, তাড়না নয়।}

\medskip\hrule\medskip

\subsection*{৩. পূর্ণ ঘটনার বিশুদ্ধ বর্ণনা (\lat{The} \lat{Full} \lat{Narration})}

আবু হুরাইরাহ (রাঃ) বলেন:

\begin{quote}
এক ব্যক্তি রাসূলুল্লাহর (সা.) কাছে এসে বলল: \textbf{\lat{“}আমি ধ্বংস হয়ে গেছি, হে আল্লাহর রাসূল!\lat{”}} নবীজি (সা.) জিজ্ঞেস করলেন: \textbf{\lat{“}তোমাকে কী ধ্বংস করল?\lat{”}} সে বলল: \lat{“}আমি রমজানে (রোজা অবস্থায়) আমার স্ত্রীর সাথে সহবাস করে ফেলেছি।\lat{”} নবীজি (সা.) জিজ্ঞেস করলেন: \textbf{\lat{“}তুমি কি একটি দাস মুক্ত করার সামর্থ্য রাখো?\lat{”}} সে বলল: \lat{“}না।\lat{”} তিনি বললেন: \textbf{\lat{“}তাহলে টানা দুই মাস রোজা রাখতে পারবে কি?\lat{”}} সে বলল: \lat{“}না।\lat{”} তিনি জিজ্ঞেস করলেন: \textbf{\lat{“}তাহলে ষাটজন মিসকিনকে খাবার দিতে পারবে কি?\lat{”}} সে বলল: \lat{“}না।\lat{”}
\end{quote}

\begin{quote}
এরপর সেই ব্যক্তি বসে রইল। এদিকে রাসূলুল্লাহর (সা.) কাছে খেজুর ভর্তি একটি বড় ঝুড়ি (\arab{عرق}) আনা হলো। নবীজি (সা.) বললেন: \textbf{\lat{“}এটা নাও এবং সদকা করে দাও।\lat{”}} সে বলল: \lat{“}আমার চেয়ে বেশি অভাবীকে দেব? মদিনার দুই পাথুরে প্রান্তরের মাঝে আমার পরিবারের চেয়ে বেশি অভাবী কোনো পরিবার নেই!\lat{”} নবীজি (সা.) এতে \textbf{হেসে ফেললেন — এমনকি তার দাঁত (মোলার) দেখা গেল} — তারপর বললেন: \textbf{\lat{“}তাহলে নাও, তোমার পরিবারকেই খাওয়াও।\lat{”}}
\end{quote}

(সহীহ বুখারি ১৯৩৬; সহীহ মুসলিম ১১১১)

\medskip\hrule\medskip

\subsection*{৪. শব্দের অর্থ (\lat{Key} \lat{Words})}

\begin{table}[h]\centering\small
\begin{tabular}{ll}
শব্দ & অর্থ \\
\hline
\textbf{হালাকতু (\arab{هلكت})} & আমি ধ্বংস হয়ে গেছি — ভেঙে পড়া অবস্থার প্রকাশ \\
\textbf{উতিকু রাকাবাহ (\arab{عتق} \arab{رقبة})} & একটি দাস/ক্রীতদাস মুক্ত করা — কাফফারার প্রথম ধাপ \\
\textbf{আল-কাফফারাহ (\arab{الكفارة})} & প্রায়শ্চিত্ত — পাপের প্রভাব ঢেকে/সংশোধনকারী আমল \\
\textbf{আল-আরাক (\arab{العرق})} & খেজুর ভর্তি বড় ঝুড়ি — আনুমানিক ১৫ সা' পরিমাণ \\
\hline
\end{tabular}
\end{table}

\medskip\hrule\medskip

\subsection*{৫. সংশ্লিষ্ট দলিল ও রেফারেন্স}

\begin{itemize}
\item \textbf{কাফফারার ধাপ (ফিকহ):} দাস মুক্ত করা $\rightarrow$ সামর্থ্য না থাকলে দুই মাস টানা রোজা $\rightarrow$ না পারলে ৬০ জন মিসকিনকে খাওয়ানো (বুখারি ১৯৩৬-এর ভিত্তিতে আলেমদের সিদ্ধান্ত)। এই ব্যক্তির ক্ষেত্রে তিনটিই সম্ভব না হওয়ায় নবীজি (সা.) খেজুর দিয়ে সদকা করালেন — পাপ থেকে বের হওয়ার পথ আল্লাহর রাসূলই দেখিয়ে দেন।
\item \textbf{আয়াত (আল-বাকারা ২:১৮৭):} \lat{“}রোজার রাতে স্ত্রীদের সাথে মেলামেশা করা তোমাদের জন্য হালাল করা হলো...\lat{”} — রাতের অনুমতির বিপরীতে দিনে নিষেধ; এই সাহাবির পাপ দিনে হওয়ায় কাফফারা প্রযোজ্য।
\item \textbf{আয়াত (আয-যুমার ৩৯:৫৩):} পাপের পর হতাশ না হওয়া — এই সাহাবি হতাশ হয়ে নিজেকে ছেড়ে দেননি; এসে পথ চেয়েছেন।
\item \textbf{হাদিস (ইবনে মাজাহ ৪২৫১):} \lat{“}প্রতিটি আদম সন্তান গুনাহ করে, সেরা গুনাহকারী তাওবাকারী।\lat{”}
\end{itemize}
\subsubsection*{মূল শিক্ষা (দলিলভিত্তিক)}

\begin{enumerate}
\item \textbf{পাপ লুকানো না — স্বীকার করা:} এই সাহাবি ভেঙে পড়ে স্বীকার করলেন; নবীজি (সা.) তাকে সমাজ থেকে বিচ্ছিন্ন করলেন না, বরং পথ দেখালেন। পাপীকে নবীজির (সা.) দরবার থেকে ফিরিয়ে দেওয়ার কোনো নজির নেই — বরং সংশোধনের পথই দেখানো হয়েছে।
\item \textbf{পাপের পর সংশোধন (\lat{remedy}):} তাওবা মানে শুধু আফসোস নয় — পাপের ক্ষতিপূরণের আমলও (কাফফারা)। নবীজি (সা.) প্রতিটি ধাপে সামর্থ্য জিজ্ঞেস করলেন — ইসলামের বিধান মানুষের ক্ষমতা অনুযায়ী (ব্যবস্থাপনা)।
\item \textbf{নবীজির (সা.) হাসি ও স্নেহ:} সবশেষে হাসিমুখে \lat{“}তোমার পরিবারকেই খাওয়াও\lat{”} — পাপ স্বীকারকারীকে নবীজি (সা.) সান্ত্বনা দিলেন; এতে পাপীর মনে আবার দাঁড়ানোর সাহস (\lat{courage}) ফিরে আসে।
\item \textbf{আল্লাহর রহমতের পথ খোলা:} এই ঘটনাটি ৩৯:৫৩-এর বাস্তব রূপ — পাপ যত বড়ই হোক, ফেরার পথ খোলা থাকলে আল্লাহর রাসূলই সেই পথ দেখান।
\end{enumerate}
\medskip\hrule\medskip

\subsection*{৬. রেফারেন্স}

\begin{itemize}
\item সহীহ বুখারি ১৯৩৬ (কিতাবুস সাওম)
\item সহীহ মুসলিম ১১১১ (কিতাবুস সিয়াম)
\item কুরআন: আল-বাকারা ২:১৮৭; আয-যুমার ৩৯:৫৩
\item ফিকহ: কাফফারার বিস্তারিত — নবীজির (সা.) এই হাদিসের ভিত্তিতে
\end{itemize}

\end{document}
